\Chapter{38}

\Verse{1}
{
    \zA{话说宝钗湘云计议已定，一宿无话。 次日湘云便请贾母等赏桂花。 贾母等都说 道：“倒是他有兴头，须要扰他这雅兴。” 至午，果然贾母带了王夫人、凤姐，兼
        请薛姨妈等进园来。 贾母因问：“那一处好？” 王夫人道：“凭老太太爱在那一处， 就在那一处。” 凤姐道：“藕香榭已经摆下了。
        那山坡下两棵桂花开的又好，河里 的水又碧清，坐在河当中亭子上，不敞亮吗？ 看看水，眼也清亮。” 贾母听了，说： “很好。” 说着，引了众人往藕香榭来。}
}
{
    \eA{After Pao-ch'ai and Hsiang-yün, we will now explain, settled everything in
        their deliberations, nothing memorable occurred, the whole night, which
        deserves to be put on record. The next day, Hsiang-yün invited dowager lady
        Chia and her other relatives to come and look at the olea flowers.}
}
{
}

\Verse{2}
{
    \zA{原来这藕香榭盖在池中，四面有窗，左右 有回廊，也是跨水接峰，后面又有曲折桥。 众人上了竹桥，凤姐忙上来搀着贾母，
        口里说道：“老祖宗只管迈大步走，不相干，这竹子桥规矩是硌吱硌吱的。” 一时进入榭中，只见栏杆外另放着两张竹案，一个上面设着杯箸酒具，一个上
        头设着茶筅茶具各色盏碟。 那边有两三个丫头煽风炉煮茶，这边另有几个丫头也煽 风炉烫酒呢。 贾母忙笑问：“这茶想的很好，且是地方东西都干净。” 湘云笑道：
        “这是宝姐姐帮着我预备的。”}
}
{
    \eA{Old lady Chia and every one else answered that as she had had the kind
        attention to ask them, they felt it their duty to avail themselves of her
        gracious invitation, much though they would be putting her to trouble and
        inconvenience.}
}
{
}

\Verse{3}
{
    \zA{贾母道：“我说那孩子细致，凡事想的妥当。” 一 面说，一面又看见柱子上挂的墨漆嵌蚌的对子，命湘云念道： 芙蓉影破归兰桨， 菱藕香深泻竹桥。
        贾母听了，又抬头看匾，因回头向薛姨妈道：“我先小时，家里也有这么一个亭子， 叫做什么枕霞阁。 我那时也只像他姐妹们这么大年纪，同着几个人，天天玩去。 谁
        知那日一下子失了脚掉下去，几乎没淹死，好容易救上来了，到底叫那木钉把头碰 破了。 如今这鬓角上那指头顶儿大的一个坑儿，就是那碰破的。}
}
{
    \eA{At twelve o'clock, therefore, old lady Chia actually took with her Madame
        Wang and lady Feng, as well as Mrs. Hsüeh and other members of her family
        whom she had asked to join them, and repaired into the garden. "Which is the
        best spot?" old lady Chia inquired. "We are ready to go wherever you may
        like, dear senior," Madame Wang ventured in response.}
}
{
}

\Verse{4}
{
    \zA{众人都怕经了水， 冒了风，说了不得了，谁知竟好了。” 凤姐不等人说，先笑道：“那时要活不得， 如今这么大福可叫谁享呢？
        可知老祖宗从小儿福寿就不小，神差鬼使，碰出那个坑 儿来，好盛福寿啊。 寿星老儿头上原是个坑儿，因为万福万寿盛满了，所以倒凸出 些来了。”
        未及说完，贾母和众人都笑软了。 贾母笑道：“这猴儿惯的了不得了， 拿着我也取起笑儿来了!恨的我撕你那油嘴。”}
}
{
    \eA{"A collation has already been spread in the Lotus Fragrance Arbour," lady
        Feng interposed. "Besides, the two olea plants, on that hill, yonder, are
        now lovely in their full blossom, and the water of that stream is jade-like
        and pellucid, so if we sit in the pavilion in the middle of it, won't we
        enjoy an open and bright view? It will be refreshing too to our eyes to
        watch the pool." "Quite right!"}
}
{
}

\Verse{5}
{
    \zA{凤姐道：“回来吃螃蟹，怕存住冷 在心里，怄老祖宗笑笑儿，就是高兴多吃两个也无妨了。” 贾母笑道：“明日叫你
        黑家白日跟着我，我倒常笑笑儿，也不许你回屋里去。” 王夫人笑道：“老太太因 为喜欢他，才惯的这么样，还这么说，他明儿越发没理了。” 贾母笑道：“我倒喜
        欢他这么着，况且他又不是那真不知高低的孩子。 家常没人，娘儿们原该说说笑笑， 横竖大礼不错就罢了。 没的倒叫他们神鬼似的做什么！” 说着，一齐进了亭子。
        献过茶，凤姐忙安放杯箸。}
}
{
    \eA{assented dowager lady Chia at this suggestion; and while expressing her
        approbation, she ushered her train of followers into the Arbour of Lotus
        Fragrance. This Arbour of Lotus Fragrance had, in fact, been erected in the
        centre of the pool. It had windows on all four sides. On the left and on the
        right, stood covered passages, which spanned the stream and connected with
        the hills.}
}
{
}

\Verse{6}
{
    \zA{上面一桌，贾母、薛姨妈、 宝钗、黛玉、宝玉； 东边一桌，湘云、王夫人、迎、探、惜； 西边靠门一小桌，李
        纨和凤姐，虚设坐位，二人皆不敢坐，只在贾母王夫人两桌上伺候。 凤姐吩咐：“螃 蟹不可多拿来，仍旧放在蒸笼里，拿十个来，吃了再拿。” 一面又要水洗了手，站
        在贾母跟前剥蟹肉。 头次让薛姨妈，薛姨妈道：“我自己掰着吃香甜，不用人让。” 凤姐便奉与贾母。 二次的便与宝玉。 又说：“把酒烫得滚热的拿来。”}
}
{
    \eA{At the back, figured a winding bridge. As the party ascended the bamboo
        bridge, lady Feng promptly advanced and supported dowager lady Chia.
        "Venerable ancestor," she said, "just walk boldly and with confident step;
        there's nothing to fear; it's the way of these bamboo bridges to go on
        creaking like this." Presently, they entered the arbour. Here they saw two
        additional bamboo tables, placed beyond the balustrade.}
}
{
}

\Verse{7}
{
    \zA{又命小丫头 们去取菊花叶儿桂花蕊熏的绿豆面子，预备着洗手。 湘云陪着吃了一个，便下座来 让人，又出至外头，命人盛两盘子给赵姨娘送去。
        又见凤姐走来道：“你张罗不惯， 你吃你的去，我先替你张罗，等散了我再吃。” 湘云不肯，又命人在那边廊上摆了 两席，让鸳鸯、琥珀、彩霞、彩云、平儿去坐。
        鸳鸯因向凤姐笑道：“二奶奶在这 里伺候，我可吃去了。” 凤姐儿道：“你们只管去，都交给我就是了。” 说着，湘 云仍入了席。 凤姐和李纨也胡乱应了个景儿。}
}
{
    \eA{On the one, were arranged cups, chopsticks and every article necessary for
        drinking wine. On the other, were laid bamboo utensils for tea, a
        tea-service and various cups and saucers. On the off side, two or three
        waiting-maids were engaged in fanning the stove to boil the water for tea.}
}
{
}

\Verse{8}
{
    \zA{凤姐仍旧下来张罗。 一时出至廊上，鸳鸯等正吃得高兴，见他来了，鸳鸯等站 起来道：“奶奶又出来做什么？ 让我们也受用一会子！” 凤姐笑道：“鸳鸯丫头越
        发坏了!我替你当差，倒不领情，还抱怨我，还不快斟一钟酒来我喝呢。” 鸳鸯笑 着，忙斟了一杯酒，送至凤姐唇边，凤姐一挺脖子喝了。 琥珀、彩霞二人也斟上一
        杯送至凤姐唇边，那凤姐也吃了。 平儿早剔了一壳黄子送来，凤姐道：“多倒些姜 醋。” 一回也吃了，笑道：“你们坐着吃罢，我可去了。”}
}
{
    \eA{On the near side were visible several other girls, who were trying with
        their fans to get a fire to light in the stove so as to warm the wines. "It
        was a capital idea," dowager lady Chia hastily exclaimed laughingly with
        vehemence, "to bring tea here. What's more, the spot and the appurtenances
        are alike so spick and span!"}
}
{
}

\Verse{9}
{
    \zA{鸳鸯笑道：“好没脸! 吃我们的东西！” 凤姐儿笑道：“你少和我作怪。 你知道你琏二爷爱上了你，要和 老太太讨了你做小老婆呢。”
        鸳鸯红了脸，咂着嘴，点着头道：“哎，这也是做奶 奶说出来的话!我不拿腥手抹你一脸算不得！” 说着站起来就要抹。 凤姐道：“好 姐姐!饶我这遭儿罢！”
        琥珀笑道：“鸳丫头要去了，平丫头还饶他？ 你们看看，他 没吃两个螃蟹，倒喝了一碟子醋了！”}
}
{
    \eA{"These things were brought by cousin Pao-ch'ai," Hsiang-yün smilingly
        explained, "so I got them ready." "This child is, I say, so scrupulously
        particular," old lady Chia observed, "that everything she does is thoroughly
        devised."}
}
{
}

\Verse{10}
{
    \zA{平儿手里正剥了个满黄螃蟹，听如此奚落他， 便拿着螃蟹照琥珀脸上来抹，口内笑骂：“我把你这嚼舌根的小蹄子儿……”琥珀 也笑着往傍边一躲。
        平儿使空了，往前一撞，恰恰的抹在凤姐腮上。 凤姐正和鸳鸯 嘲笑，不防吓了一跳，“嗳哟”了一声，众人掌不住都哈哈的大笑起来。 凤姐也禁
        不住笑骂道：“死娼妇!吃离了眼了!混抹你娘的！” 平儿忙赶过来替他擦了，亲自 去端水。 鸳鸯道：“阿弥陀佛!这才是现报呢。”}
}
{
    \eA{As she gave utterance to her feelings, her attention was attracted by a pair
        of scrolls of black lacquer, inlaid with mother-of-pearl, suspended on the
        pillars, and she asked Hsiang-yün to tell her what the mottoes were. The
        text she read was:  Snapped is the shade of the hibiscus by the fragrant oar
        of a boat  homeward bound. Deep flows the perfume of the lily and the lotus
        underneath the  bamboo  bridge.}
}
{
}

\Verse{11}
{
    \zA{贾母那边听见，一叠连声问：“见 了什么了，这么乐？ 告诉我们也笑笑。” 鸳鸯等忙高声笑回道：“二奶奶来抢螃蟹
        吃，平儿恼了，抹了他主子一脸螃蟹黄子：主子奴才打架呢！” 贾母和王夫人等听 了，也笑起来。
        贾母笑道：“你们看他可怜见儿的，那小腿子、脐子给他点子吃罢。” 鸳鸯等笑着答应了，高声的说道：“这满桌子的腿子，二奶奶只管吃就是了。” 凤
        姐笑着洗了脸，走来又伏侍贾母等吃了一回。 黛玉弱不敢多吃，只吃了一点夹子肉就下来了。}
}
{
    \eA{After listening to the motto, old lady Chia raised her head and cast a
        glance upon the tablet; then turning round: "Long ago, when I was young,"
        she observed, addressing herself to Mrs. Hsüeh, "we likewise had at home a
        pavilion like this called 'the Hall reclining on the russet clouds,' or some
        other such name.}
}
{
}

\Verse{12}
{
    \zA{贾母一时也不吃了。 大家都洗 了手。 也有看花的，也有弄水看鱼的，游玩了一回。 王夫人因问贾母：“这里风大， 才又吃了螃蟹，老太太还是回屋里去歇歇罢。
        若高兴，明日再来逛逛。” 贾母听了， 笑道：“正是呢。 我怕你们高兴，我走了，又怕扫了你们的兴； 既这么说，咱们就 都去罢。”
        回头嘱咐湘云：“别让你宝哥哥多吃了。” 湘云答应着。 又嘱咐湘云、 宝钗二人说：“你们两个也别多吃了。 那东西虽好吃，不是什么好的，吃多了肚子 疼。”}
}
{
    \eA{At that time, I was of the same age as the girls, and my wont was to go day
        after day and play with my sisters there. One day, I, unexpectedly, slipped
        and fell into the water, and I had a narrow escape from being drowned; for
        it was after great difficulty, that they managed to drag me out safe and
        sound. But my head was, after all, bumped about against the wooden nails;}
}
{
}

\Verse{13}
{
    \zA{二人忙应着。 送出园外，仍旧回来，命将残席收拾了另摆。 宝玉道：“也不 用摆，咱们且做诗。 把那大团圆桌子放在当中，酒菜都放着。 也不必拘定坐位，有
        爱吃的去吃，大家散坐，岂不便宜？” 宝钗道：“这话极是。” 湘云道：“虽这么 说，还有别人。” 因又命另摆一桌，拣了热螃蟹来，请袭人、紫鹃、司棋、侍书、
        入画、莺儿、翠墨等一处共坐。 山坡桂树底下铺下两条花毯，命支应的婆子并小丫 头等也都坐了，只管随意吃喝，等使唤再来。}
}
{
    \eA{so much so, that this hole of the length of a finger, which you can see up
        to this day on my temple, comes from the bruises I sustained. All my people
        were in a funk that I'd be the worse for this ducking and continued in fear
        and trembling lest I should catch a chill. 'It was dreadful, dreadful!' they
        opined, but I managed, little though every one thought it, to keep in
        splendid health."}
}
{
}

\Verse{14}
{
    \zA{湘云便取了诗题，用针绾在墙上。 众人看了，都说：“新奇!只怕做不出来。” 湘云又把不限韵的缘故说了一番，宝玉道：“这才是正理。 我也最不喜限韵。” 黛
        玉因不大吃酒，又不吃螃蟹，自命人掇了一个绣墩，倚栏坐着，拿着钓杆钓鱼。 宝 钗手里拿着一枝桂花，玩了一回，俯在窗槛上，掐了桂蕊，扔在水面，引的那游鱼 ？
        上来唼喋。 湘云出一回神，又让一回袭人等，又招呼山坡下的众人只管放量吃。 探春和李纨、惜春正立在垂柳阴中看鸥鹭。}
}
{
    \eA{Lady Feng allowed no time to any one else to put in a word; but anticipating
        them: "Had you then not survived, who would now be enjoying these immense
        blessings!" she smiled. "This makes it evident that no small amount of
        happiness and long life were in store for you, venerable ancestor, from your
        very youth up!}
}
{
}

\Verse{15}
{
    \zA{迎春却独在花阴下，拿着个针儿穿茉莉 花。 宝玉又看了一回黛玉钓鱼，一回又俯在宝钗傍边说笑两句，一回又看袭人等吃
        螃蟹，自己也陪他喝两口酒，袭人又剥一壳肉给他吃。 黛玉放下钓杆，走至座间，拿起那乌梅银花自斟壶来，拣了一个小小的海棠冻 石蕉叶杯。
        丫头看见，知他要饮酒，忙着走上来斟。 黛玉道：“你们只管吃去，让 我自己斟才有趣儿。” 说着便斟了半盏看时，却是黄酒，因说道：“我吃了一点子
        螃蟹，觉得心口微微的疼，须得热热的吃口烧酒。”}
}
{
    \eA{It was by the agency of the spirits that this hole was knocked open so that
        they might fill it up with happiness and longevity! The old man Shou Hsing
        had, in fact, a hole in his head, which was so full of every kind of
        blessing conducive to happiness and long life that it bulged up ever so
        high!"}
}
{
}

\Verse{16}
{
    \zA{宝玉忙接道：“有烧酒。” 便 命将那合欢花浸的酒烫一壶来，黛玉也只吃了一口便放下了。 宝钗也走过来，另拿
        了一只杯来，也饮了一口放下，便蘸笔至墙上把头一个《忆菊》勾了，底下又赘一 个“蘅”字。 宝玉忙道：“好姐姐，第二个我已有了四句了，你让我做罢。” 宝钗
        笑道：“我好容易有了一首，你就忙的这样。” 黛玉也不说话，接过笔来把第八个 《问菊》勾了，接着把第十一个《菊梦》也勾了，也赘上了一个“潇”字。}
}
{
    \eA{Before, however, she could conclude, dowager lady Chia and the rest were
        convulsed with such laughter that their bodies doubled in two. "This monkey
        is given to dreadful tricks!" laughed old lady Chia. "She's always ready to
        make a scapegoat of me to evoke amusement. But would that I could take that
        glib mouth of yours and rend it in pieces."}
}
{
}

\Verse{17}
{
    \zA{宝玉也 拿起笔来将第二个《访菊》也勾了，也赘上一个“怡”字。 探春起来看着道：“竟 没人作《簪菊》？ 让我作。”
        又指着宝玉笑道：“才宣过：总不许带出闺阁字样来， 你可要留神。” 说着，只见湘云走来，将第四第五《对菊》《供菊》一连两个都勾 了，也赘上一个“湘”字。
        探春道：“你也该起个号。” 湘云笑道：“我们家里如 今虽有几处轩馆，我又不住着，借了来也没趣。” 宝钗笑道：“方才老太太说，你
        们家里也有一个水亭，叫做枕霞阁，难道不是你的？}
}
{
    \eA{"It's because I feared that the cold might, when you by and bye have some
        crabs to eat, accumulate in your intestines," lady Feng pleaded, "that I
        tried to induce you, dear senior, to have a laugh, so as to make you gay and
        merry. For one can, when in high spirits, indulge in a couple of them more
        with impunity."}
}
{
}

\Verse{18}
{
    \zA{如今虽没了，你到底是旧主人。” 众人都道：“有理。” 宝玉不待湘云动手，便代将“湘”字抹了，改了一个“霞” 字。
        没有顿饭工夫，十二题已全，各自誊出来，都交与迎春，另拿了一张雪浪笺过 来，一并誊录出来。 某人作的底下赘明某人的号。 李纨等从头看道： 忆　菊　　蘅芜君
        怅望西风抱闷思，蓼红苇白断肠时。 空篱旧圃秋无迹，冷月清霜梦有知。 念念心随归雁远，寥寥坐听晚砧迟。 谁怜我为黄花瘦，慰语重阳会有期。}
}
{
    \eA{"By and bye," smiled old lady Chia, "I'll make you follow me day and night,
        so that I may constantly be amused and feel my mind diverted; I won't let
        you go back to your home."}
}
{
}

\Verse{19}
{
    \zA{访　菊　　怡红公子 闲趁霜晴试一游，酒杯药盏莫淹留。 霜前月下谁家种？ 槛外篱边何处秋？ 蜡屐远来情得得，冷吟不尽兴悠悠。
        黄花若解怜诗客，休负今朝挂杖头。 种　菊　　怡红公子 携锄秋圃自移来，篱畔庭前处处栽。 昨夜不期经雨活，今朝犹喜带霜开。
        冷吟秋色诗千首，醉酹寒香酒一杯。 泉溉泥封勤护惜，好和井径绝尘埃。 对　菊　　枕霞旧友 别圃移来贵比金，一丛浅淡一丛深。
        萧疏篱畔科头坐，清冷香中抱膝吟。}
}
{
    \eA{"It's that weakness of yours for her, venerable senior," Madame Wang
        observed with a smile, "that has got her into the way of behaving in this
        manner, and, if you go on speaking to her as you do, she'll soon become ever
        so much the more unreasonable." "I like her such as she is," dowager lady
        Chia laughed. "Besides, she's truly no child, ignorant of the distinction
        between high and low.}
}
{
}

\Verse{20}
{
    \zA{数去更无君傲世，看来惟有我知音! 秋光荏苒休孤负，相对原宜惜寸阴。 供　菊　　枕霞旧友 弹琴酌酒喜堪俦，几案婷婷点缀幽。
        隔坐香分三径露，抛书人对一枝秋。 霜清纸帐来新梦，圃冷斜阳忆旧游。 傲世也因同气味，春风桃李未淹留。 咏　菊　　潇湘妃子
        无赖诗魔昏晓侵，绕篱欹石自沉音。 毫端蕴秀临霜写，口角噙香对月吟。 满纸自怜题素怨，片言谁解诉秋心？ 一从陶令评章后，千古高风说到今。 画　菊　　蘅芜君
        诗馀戏笔不知狂，岂是丹青费较量？}
}
{
    \eA{When we are at home, with no strangers present, we ladies should be on terms
        like these, and as long, in fact, as we don't overstep propriety, it's all
        right. If not, what would be the earthly use of making them behave like so
        many saints?" While bandying words, they entered the pavilion in a body.
        After tea, lady Feng hastened to lay out the cups and chopsticks.}
}
{
}

\Verse{21}
{
    \zA{聚叶泼成千点墨，攒花染出几痕霜。 淡浓神会风前影，跳脱秋生腕底香。 莫认东篱闲采掇，粘屏聊以慰重阳。 问　菊　　潇湘妃子
        欲讯秋情众莫知，喃喃负手扣东篱。 孤标傲世偕谁隐？ 一样开花为底迟？ 圃露庭霜何寂寞？ 雁归蛩病可相思？ 莫言举世无谈者，解语何妨话片时？
        簪　菊　　蕉下客 瓶供篱栽日日忙，折来休认镜中妆。 长安公子因花癖，彭泽先生是酒狂。 短鬓冷沾三径露，葛巾香染九秋霜。 高情不入时人眼，拍手凭他笑路旁。}
}
{
    \eA{At the upper table then seated herself old lady Chia, Mrs. Hsüeh, Pao-ch'ai,
        Tai-yü and Pao-yü. Round the table, on the east, sat Shih Hsiang-yün, Madame
        Wang, Ying Ch'un, T'an Ch'un and Hsi Ch'un. At the small table, leaning
        against the door on the west side, Li Wan and lady Feng assigned themselves
        places.}
}
{
}

\Verse{22}
{
    \zA{菊　影　　枕霞旧友 秋光叠叠复重重，潜度偷移三径中。 窗隔疏灯描远近，篱筛破月锁玲珑。 寒芳留照魂应驻，霜印传神梦也空。
        珍重暗香踏碎处，凭谁醉眼认朦胧。 菊　梦　　潇湘妃子 篱畔秋酣一觉清，和云伴月不分明。 登仙非慕庄生蝶，忆旧还寻陶令盟。
        睡去依依随雁断，惊回故故恼蛩鸣。 醒时幽怨同谁诉：衰草寒烟无限情! 残　菊　　蕉下客 露凝霜重渐倾欹，宴赏才过小雪时。 蒂有馀香金淡泊，枝无全叶翠离披。
        半床落月蛩声切，万里寒云雁阵迟。}
}
{
    \eA{But it was for the mere sake of appearances, as neither of them ventured to
        sit down, but remained in attendance at the two tables, occupied by old lady
        Chia and Madame Wang. "You'd better," lady Feng said, "not bring in too many
        crabs at a time. Throw these again into the steaming-basket! Only serve ten;
        and when they're eaten, a fresh supply can be fetched!"}
}
{
}

\Verse{23}
{
    \zA{明岁秋分知再会，暂时分手莫相思! 众人看一首，赞一首，彼此称扬不绝。 李纨笑道：“等我从公评来。 通篇看来， 各人有各人的警句。
        今日公评：《咏菊》第一，《问菊》第二，《菊梦》第三，题 目新，诗也新，立意更新了，只得要推潇湘妃子为魁了。 然后《簪菊》、《对菊》、
        《供菊》、《画菊》、《忆菊》次之。” 宝玉听说，喜的拍手叫道：“极是!极公！” 黛玉道：“我那个也不好，到底伤于纤巧些。”
        李纨道：“巧的却好，不露堆砌生 硬。”}
}
{
    \eA{Asking, at the same time, for water, she washed her hands, and, taking her
        position near dowager lady Chia, she scooped out the meat from a crab, and
        offered the first help to Mrs. Hsüeh. "They'll be sweeter were I to open
        them with my own hands," Mrs. Hsüeh remarked, "there's no need for any one
        to serve me." Lady Feng, therefore, presented it to old lady Chia and handed
        a second portion to Pao-yü.}
}
{
}

\Verse{24}
{
    \zA{黛玉道：“据我看来，头一句好的是‘圃冷斜阳忆旧游’，这句背面傅粉； ‘抛书人对一枝秋’，已经妙绝，将供菊说完，没处再说，故翻回来想到未折未供
        之先，意思深远！” 李纨笑道：“固如此说，你的‘口角噙香’一句也敌得过了。” 探春又道：“到底要算蘅芜君沉着，‘秋无迹’，‘梦有知’，把个‘忆’字竟烘
        染出来了。” 宝钗笑道：“你的‘短鬓冷沾’，‘葛巾香染’，也就把簪菊形容的 一个缝儿也没有。”}
}
{
    \eA{"Make the wine as warm as possible and bring it in!" she then went on to
        cry. "Go," she added, directing the servant-girls, "and fetch the powder,
        made of green beans, and scented with the leaves of chrysanthemums and the
        stamens of the olea fragrans; and keep it ready to rinse our hands with."}
}
{
}

\Verse{25}
{
    \zA{湘云笑道：“‘偕谁隐’，‘为底迟’，真真把个菊花问的无 言可对！” 李纨笑道：“那么着，像‘科头坐’，‘抱膝吟’，竟一时也舍不得离
        了菊花，菊花有知，倒还怕腻烦了呢！” 说的大家都笑了。 宝玉笑道：“这场我又 落第了。 难道‘谁家种’，‘何处秋’，‘蜡屐远来’，‘冷吟不尽’，那都不是
        访不成？ ‘昨夜雨’，‘今朝霜’，都不是种不成？}
}
{
    \eA{Shih Hsiang-yün had a crab to bear the others company, but no sooner had she
        done than she retired to a lower seat, from where she helped her guests.
        When she, however, walked out a second time to give orders to fill two
        dishes and send them over to Mrs. Chao, she perceived lady Feng come up to
        her again. "You're not accustomed to entertaining," she said, "so go and
        have your share to eat.}
}
{
}

\Verse{26}
{
    \zA{但恨敌不上‘口角噙香对月吟’、 ‘清冷香中抱膝吟’、‘短鬓’、‘葛巾’、‘金淡泊’、‘翠离披’、‘秋无迹’、 ‘梦有知’这几句罢了。”
        又道：“明日闲了，我一个人做出十二首来。” 李纨道： “你的也好，只是不及这几句新雅就是了。” 大家又评了一回，复又要了热螃蟹来，就在大圆桌上吃了一回。
        宝玉笑道：“今 日持螯赏桂，亦不可无诗，我已吟成，谁还敢作？” 说着，便忙洗了手，提笔写出， 众人看道： 持螯更喜桂阴凉，泼醋擂姜兴欲狂。}
}
{
    \eA{I'll attend to the people for you first, and, when they've gone, I'll have
        all I want." Hsiang-yün would not agree to her proposal. But giving further
        directions to the servants to spread two tables under the verandah on the
        off-side, she pressed Yüan Yang, Hu Po, Ts'ai Hsia, Ts'ai Yün and P'ing Erh
        to go and seat themselves.}
}
{
}

\Verse{27}
{
    \zA{饕餮王孙应有酒，横行公子竟无肠! 脐间积冷馋忘忌，指上沾腥洗尚香。 原为世人美口腹，坡仙曾笑一生忙。 黛玉笑道：“这样的诗，一时要一百首也有。”
        宝玉笑道：“你这会子才力已尽， 不说不能作了，还褒贬人家。” 黛玉听了，也不答言，略一仰首微吟，提起笔来一 挥，已有了一首。 众人看道：
        铁甲长戈死未忘，堆盘色相喜先尝。 螯封嫩玉双双满，壳凸红脂块块香。 多肉更怜卿八足，助情谁劝我千觞？ 对兹佳品酬佳节，桂拂清风菊带霜。}
}
{
    \eA{"Lady Secunda," consequently ventured Yüan Yang, "you're in here doing the
        honours, so may I go and have something to eat?" "You can all go," replied
        lady Feng; "leave everything in my charge, and it will be all right." While
        these words were being spoken, Shih Hsiang-yün resumed her place at the
        banquet. Lady Feng and Li Wan then took hurry-scurry something to eat as a
        matter of form;}
}
{
}

\Verse{28}
{
    \zA{宝玉看了，正喝彩时，黛玉便一把撕了，命人烧去，因笑道：“我做的不及你的， 我烧了罢。 你那个很好，比方才的菊花诗还好，你留着他给人看看。”
        宝钗笑道：“我也勉强了一首，未必好，写出来取笑儿罢。” 说着，也写出来。 大家看时，写道： 桂霭桐阴坐举觞，长安涎口盼重阳。
        眼前道路无经纬，皮里春秋空黑黄。 看到这里，众人不禁叫绝。 宝玉道：“骂得痛快!我的诗也该烧了。” 看底下道： 酒未涤腥还用菊，性防积冷定须姜。
        于今落釜成何益？}
}
{
    \eA{but lady Feng came down once more to look after things. After a time, she
        stepped out on the verandah where Yüan Yang and the other girls were having
        their refreshments in high glee. As soon as they caught sight of her, Yuan
        Yang and her companions stood up. "What has your ladyship come out again
        for?" they inquired. "Do let us also enjoy a little peace and quiet!" "This
        chit Yüan Yang is worse than ever!"}
}
{
}

\Verse{29}
{
    \zA{月浦空馀禾黍香。 众人看毕，都说：“这方是食蟹的绝唱!这些小题目，原要寓大意思，才算是大才。 只是讽刺世人太毒了些。” 说着，只见平儿复进园来。
        不知却做什么，且听下回分解。 (本章完)}
}
{
    \eA{lady Feng laughed. "Here I'm slaving away for you, and, instead of feeling
        grateful to me, you bear me a grudge! But don't you yet quick pour me a cup
        of wine?" Yüan Yang immediately smiled, and filling a cup, she applied it to
        lady Feng's lips. Lady Feng stretched out her neck and emptied it. But Hu Po
        and Ts'ai Hsia thereupon likewise replenished a cup and put it to lady
        Feng's mouth.}
}
{
}

\Verse{30}
{
    \zA{}
}
{
    \eA{Lady Feng swallowed the contents of that as well. P'ing Erh had, by this
        time, brought her some yellow meat which she had picked out from the shell.
        "Pour plenty of ginger and vinegar!" shouted lady Feng, and, in a moment,
        she made short work of that too. "You people," she smiled, "had better sit
        down and have something to eat, for I'm off now." "You brazen-faced thing,"
        exclaimed Yüan Yang laughingly, "to eat what was intended for us!" "Don't be
        so captious with me!" smiled lady Feng. "Are you aware that your master
        Secundus, Mr. Lien, has taken such a violent fancy to you that he means to
        speak to our old lady to let you be his secondary wife!" Yüan Yang blushed
        crimson. "Ts'ui!" she shouted. "Are these really words to issue from the
        mouth of a lady!}
}
{
}

\Verse{31}
{
    \zA{}
}
{
    \eA{But if I don't daub your face all over with my filthy hands, I won't feel
        happy!" Saying this, she rushed up to her. She was about to besmear her
        face, when lady Feng pleaded: "My dear child, do let me off this time!" "Lo,
        that girl Yüan," laughed Hu Po, "wishes to smear her, and that hussey P'ing
        still spares her! Look here, she has scarcely had two crabs, and she has
        drunk a whole saucerful of vinegar!" P'ing Erh was holding a crab full of
        yellow meat, which she was in the act of cleaning. As soon therefore as she
        heard this taunt, she came, crab in hand, to spatter Hu Po's face, as she
        laughingly reviled her. "I'll take you minx with that cajoling tongue of
        yours" she cried, "and…." But, Hu Po, while also indulging in laughter, drew
        aside;}
}
{
}

\Verse{32}
{
    \zA{}
}
{
    \eA{so P'ing Erh beat the air, and fell forward, daubing, by a strange
        coincidence, the cheek of lady Feng. Lady Feng was at the moment having a
        little good-humoured raillery with Yüan Yang, and was taken so much off her
        guard, that she was quite startled out of her senses. "Ai-yah!" she
        ejaculated. The bystanders found it difficult to keep their countenance,
        and, with one voice, they exploded into a boisterous fit of laughter. Lady
        Feng as well could not help feeling amused, and smilingly she upbraided her.
        "You stupid wench!" she said; "Have you by gorging lost your eyesight that
        you recklessly smudge your mistress' face?"}
}
{
}

\Verse{33}
{
    \zA{}
}
{
    \eA{P'ing Erh hastily crossed over and wiped her face for her, and then went in
        person to fetch some water. "O-mi-to-fu," ejaculated Yüan Yang, "this is a
        distinct retribution!" Dowager lady Chia, though seated on the other side,
        overheard their shouts, and she consecutively made inquiries as to what they
        had seen to tickled their fancy so. "Tell us," (she urged), "what it is so
        that we too should have a laugh." "Our lady Secunda," Yüan Yang and the
        other maids forthwith laughingly cried, "came to steal our crabs and eat
        them, and P'ing Erh got angry and daubed her mistress' face all over with
        yellow meat. So our mistress and that slave-girl are now having a scuffle
        over it." This report filled dowager lady Chia, Madame Wang and the other
        inmates with them with much merriment.}
}
{
}

\Verse{34}
{
    \zA{}
}
{
    \eA{"Do have pity on her," dowager lady Chia laughed, "and let her have some of
        those small legs and entrails to eat, and have done!" Yuan Yang and her
        companions assented, much amused. "Mistress Secunda," they shouted in a loud
        tone of voice, "you're at liberty to eat this whole tableful of legs!" But
        having washed her face clean, lady Feng approached old lady Chia and the
        other guests and waited upon them for a time, while they partook of
        refreshments. Tai-yü did not, with her weak physique, venture to overload
        her stomach, so partaking of a little meat from the claws, she left the
        table. Presently, however, dowager lady Chia too abandoned all idea of
        having anything more to eat.}
}
{
}

\Verse{35}
{
    \zA{}
}
{
    \eA{The company therefore quitted the banquet; and, when they had rinsed their
        hands, some admired the flowers, some played with the water, others looked
        at the fish. After a short stroll, Madame Wang turned round and remarked to
        old lady Chia: "There's plenty of wind here. Besides, you've just had crabs;
        so it would be prudent for you, venerable senior, to return home and rest.
        And if you feel in the humour, we can come again for a turn to-morrow."
        "Quite true!" acquiesced dowager lady Chia, in reply to this suggestion. "I
        was afraid that if I left, now that you're all in exuberant spirits, I
        mightn't again be spoiling your fun, (so I didn't budge). But as the idea
        originates from yourselves do go as you please, (while I retire).}
}
{
}

\Verse{36}
{
    \zA{}
}
{
    \eA{But," she said to Hsiang-yün, "don't allow your cousin Secundus, Pao-yü, and
        your cousin Lin to have too much to eat." Then when Hsiang-yün had signified
        her obedience, "You two girls," continuing, she recommended Hsiang-yün and
        Pao-ch'ai, "must not also have more than is good for you. Those things are,
        it's true, luscious, but they're not very wholesome; and if you eat
        immoderately of them, why, you'll get stomachaches." Both girls promised
        with alacrity to be careful; and, having escorted her beyond the confines of
        the garden, they retraced their steps and ordered the servants to clear the
        remnants of the banquet and to lay out a new supply of refreshments.
        "There's no use of any regular spread out!" Pao-yü interposed.}
}
{
}

\Verse{37}
{
    \zA{}
}
{
    \eA{"When you are about to write verses, that big round table can be put in the
        centre and the wines and eatables laid on it. Neither will there be any need
        to ceremoniously have any fixed seats. Let those who may want anything to
        eat, go up to it and take what they like; and if we seat ourselves,
        scattered all over the place, won't it be far more convenient for us?" "Your
        idea is excellent!" Pao-ch'ai answered. "This is all very well," Hsiang-yün
        observed, "but there are others to be studied besides ourselves!" Issuing
        consequently further directions for another table to be laid, and picking
        out some hot crabs, she asked Hsi Jen, Tzu Chüan, Ssu Ch'i, Shih Shu, Ju
        Hua, Ying Erh, Ts'ui Mo and the other girls to sit together and form a
        party.}
}
{
}

\Verse{38}
{
    \zA{}
}
{
    \eA{Then having a couple of flowered rugs spread under the olea trees on the
        hills, she bade the matrons on duty, the waiting-maids and other servants to
        likewise make themselves comfortable and to eat and drink at their pleasure
        until they were wanted, when they could come and answer the calls.
        Hsiang-yün next fetched the themes for the verses and pinned them with a
        needle on the wall. "They're full of originality," one and all exclaimed
        after perusal, "we fear we couldn't write anything on them." Hsiang-yün then
        went onto explain to them the reasons that had prompted her not to determine
        upon any particular rhymes. "Yes, quite right!" put in Pao-yü. "I myself
        don't fancy hard and fast rhymes!"}
}
{
}

\Verse{39}
{
    \zA{}
}
{
    \eA{But Lin Tai-yü, being unable to stand much wine and to take any crabs, told,
        on her own account, a servant to fetch an embroidered cushion; and, seating
        herself in such a way as to lean against the railing, she took up a
        fishing-rod and began to fish. Pao-ch'ai played for a time with a twig of
        olea she held in her hand, then resting on the window-sill, she plucked the
        petals, and threw them into the water, attracting the fish, which went by,
        to rise to the surface and nibble at them.}
}
{
}

\Verse{40}
{
    \zA{}
}
{
    \eA{Hsiang-yün, after a few moments of abstraction, urged Hsi Jen and the other
        girls to help themselves to anything they wanted, and beckoned to the
        servants, seated at the foot of the hill, to eat to their heart's content.
        Tan Ch'un, in company with Li Wan and Hsi Ch'un, stood meanwhile under the
        shade of the weeping willows, and looked at the widgeons and egrets. Ying
        Ch'un, on the other hand, was all alone under the shade of some trees,
        threading double jasmine flowers, with a needle specially adapted for the
        purpose. Pao-yü too watched Tai-yü fishing for a while. At one time he leant
        next to Pao-ch'ai and cracked a few jokes with her. And at another, he
        drank, when he noticed Hsi Jen feasting on crabs with her companions, a few
        mouthfuls of wine to keep her company.}
}
{
}

\Verse{41}
{
    \zA{}
}
{
    \eA{At this, Hsi Jen cleaned the meat out of a shell, and gave it to him to eat.
        Tai-yü then put down the fishing-rod, and, approaching the seats, she laid
        hold of a small black tankard, ornamented with silver plum flowers, and
        selected a tiny cup, made of transparent stone, red like a begonia, and in
        the shape of a banana leaf. A servant-girl observed her movements, and,
        concluding that she felt inclined to have a drink, she drew near with
        hurried step to pour some wine for her. "You girls had better go on eating,"
        Tai-yü remonstrated, "and let me help myself; there'll be some fun in it
        then!" So speaking, she filled for herself a cup half full;}
}
{
}

\Verse{42}
{
    \zA{}
}
{
    \eA{but discovering that it was yellow wine, "I've eaten only a little bit of
        crab," she said, "and yet I feel my mouth slightly sore; so what would do
        for me now is a mouthful of very hot distilled spirit." Pao-yü hastened to
        take up her remark. "There's some distilled spirit," he chimed in. "Take
        some of that wine," he there and then shouted out to a servant, "scented
        with acacia flowers, and warm a tankard of it." When however it was brought
        Tai-yü simply took a sip and put it down again. Pao-ch'ai too then came
        forward, and picked up a double cup; but, after drinking a mouthful of it,
        she lay it aside, and, moistening her pen, she walked up to the wall, and
        marked off the first theme: "longing for chrysanthemums," below which she
        appended a character "Heng." "My dear cousin," promptly remarked Pao-yü.}
}
{
}

\Verse{43}
{
    \zA{}
}
{
    \eA{"I've already got four lines of the second theme so let me write on it!" "I
        managed, after ever so much difficulty, to put a stanza together," Pao-ch'ai
        smiled, "and are you now in such a hurry to deprive me of it?" Without so
        much as a word, Tai-yü took a pen and put a distinctive sign opposite the
        eighth, consisting of: "ask the chrysanthemums;" and, singling out, in quick
        succession, the eleventh: "dream of chrysanthemums," as well, she too
        affixed for herself the word "Hsiao" below. But Pao-yü likewise got a pen,
        and marked his choice, the twelfth on the list: "seek for chrysanthemums,"
        by the side of which he wrote the character "Chiang." T'an Ch'un thereupon
        rose to her feet.}
}
{
}

\Verse{44}
{
    \zA{}
}
{
    \eA{"If there's no one to write on 'Pinning the chrysanthemums'" she observed,
        while scrutinising the themes, "do let me have it! It has just been ruled,"
        she continued, pointing at Pao-yü with a significant smile, "that it is on
        no account permissible to introduce any expressions, bearing reference to
        the inner chambers, so you'd better be on your guard!" But as she spoke, she
        perceived Hsiang-yün come forward, and jointly mark the fourth and fifth,
        that is: "facing the chrysanthemums," and "putting chrysanthemums in vases,"
        to which she, like the others, appended a word, Hsiang." "You too should get
        a style or other!" T'an Ch'un suggested.}
}
{
}

\Verse{45}
{
    \zA{}
}
{
    \eA{"In our home," smiled Hsiang-yün, "there exist, it is true, at present
        several halls and structures, but as I don't live in either, there'll be no
        fun in it were I to borrow the name of any one of them!" "Our venerable
        senior just said," Pao-ch'ai observed laughingly, "that there was also in
        your home a water-pavilion called 'leaning on russet clouds hall,' and is it
        likely that it wasn't yours? But albeit it doesn't exist now-a-days, you
        were anyhow its mistress of old." "She's right!" one and all exclaimed.
        Pao-yü therefore allowed Hsiang-yün no time to make a move, but forthwith
        rubbed off the character "Hsiang," for her and substituted that of "Hsia"
        (russet). A short time only elapsed before the compositions on the twelve
        themes had all been completed.}
}
{
}

\Verse{46}
{
    \zA{}
}
{
    \eA{After they had each copied out their respective verses, they handed them to
        Ying Ch'un, who took a separate sheet of snow-white fancy paper, and
        transcribed them together, affixing distinctly under each stanza the style
        of the composer. Li Wan and her assistants then began to read, starting from
        the first on the list, the verses which follow: "Longing for
        chrysanthemums," by the "Princess of Heng Wu." With anguish sore I face the
        western breeze, and wrapt in grief, I  pine for you! What time the smart
        weed russet turns, and the reeds white, my heart  is rent in two. When in
        autumn the hedges thin, and gardens waste, all trace of you  is  gone. When
        the moon waxeth cold, and the dew pure, my dreams then know  something of
        you. With constant yearnings my heart follows you as far as wild geese
        homeward fly.}
}
{
}

\Verse{47}
{
    \zA{}
}
{
    \eA{Lonesome I sit and lend an ear, till a late hour to the sound of the  block!
        For you, ye yellow flowers, I've grown haggard and worn, but who doth  pity
        me,  And breathe one word of cheer that in the ninth moon I will soon meet
        you again? "Search for chrysanthemums," by the "Gentleman of I Hung:"  When
        I have naught to do, I'll seize the first fine day to try and  stroll about.
        Neither wine-cups nor cups of medicine will then deter me from my  wish. Who
        plants the flowers in all those spots, facing the dew and under  the moon's
        rays? Outside the rails they grow and by the hedge; but in autumn where do
        they go? With sandals waxed I come from distant shores; my feelings all
        exuberant; But as on this cold day I can't exhaust my song, my spirits get
        depressed.}
}
{
}

\Verse{48}
{
    \zA{}
}
{
    \eA{The yellow flowers, if they but knew how comfort to a poet to afford,  Would
        not let me this early morn trudge out in vain with my  cash-laden  staff.
        "Planting chrysanthemums," by the Gentleman of "I Hung:"  When autumn
        breaks, I take my hoe, and moving them myself out of the  park,  I plant
        them everywhere near the hedges and in the foreground of the  halls. Last
        night, when least expected, they got a good shower, which made  them all
        revive. This morn my spirits still rise high, as the buds burst in bloom
        bedecked with frost. Now that it's cool, a thousand stanzas on the autumn
        scenery I sing. In ecstasies from drink, I toast their blossom in a cup of
        cold, and  fragrant wine. With spring water.}
}
{
}

\Verse{49}
{
    \zA{}
}
{
    \eA{I sprinkle them, cover the roots with mould and  well tend them,  So that
        they may, like the path near the well, be free of every grain  of dirt.
        "Facing the chrysanthemums," by the "Old friend of the Hall reclining on the
        russet clouds." From other gardens I transplant them, and I treasure them
        like gold. One cluster bears light-coloured bloom; another bears dark
        shades. I sit with head uncovered by the sparse-leaved artemesia hedge,  And
        in their pure and cool fragrance, clasping my knees, I hum my  lays. In the
        whole world, methinks, none see the light as peerless as these  flowers.
        From all I see you have no other friend more intimate than me. Such autumn
        splendour, I must not misuse, as steadily it fleets. My gaze I fix on you as
        I am fain each moment to enjoy!}
}
{
}

\Verse{50}
{
    \zA{}
}
{
    \eA{"Putting chrysanthemums in vases," by the "Old Friend of the hall reclining
        on the russet clouds." The lute I thrum, and quaff my wine, joyful at heart
        that ye are meet  to be my mates. The various tables, on which ye are laid,
        adorn with beauteous grace  this quiet nook. The fragrant dew, next to the
        spot I sit, is far apart from that by  the three paths. I fling my book
        aside and turn my gaze upon a twig full of your  autumn  (bloom). What time
        the frost is pure, a new dream steals o'er me, as by the  paper screen I
        rest. When cold holdeth the park, and the sun's rays do slant, I long and
        yearn for you, old friends. I too differ from others in this world, for my
        own tastes resemble  those of yours.}
}
{
}

\Verse{51}
{
    \zA{}
}
{
    \eA{The vernal winds do not hinder the peach tree and the pear from  bursting
        forth in bloom. "Singing chrysanthemums," by the "Hsiao Hsiang consort."
        Eating the bread of idleness, the frenzy of poetry creeps over me  both
        night and day. Round past the hedge I wend, and, leaning on the rock, I
        intone  verses  gently to myself. From the point of my pencil emanate lines
        of recondite grace, so near  the frost I write. Some scent I hold by the
        side of my mouth, and, turning to the moon,  I  sing my sentiments. With
        self-pitying lines pages I fill, so as utterance to give to all  my cares
        and woes. From these few scanty words, who could fathom the secrets of my
        heart  about the autumntide?}
}
{
}

\Verse{52}
{
    \zA{}
}
{
    \eA{Beginning from the time when T'ao, the magistrate, did criticise the  beauty
        of your bloom,  Yea, from that date remote up to this very day, your high
        renown has  ever been extolled. "Drawing chrysanthemums," by the "Princess
        of Heng Wu." Verses I've had enough, so with my pens I play; with no idea
        that I  am  mad. Do I make use of pigments red or green as to involve a task
        of  toilsome work? To form clusters of leaves, I sprinkle simply here and
        there a  thousand specks of ink. And when I've drawn the semblance of the
        flowers, some spots I make  to  represent the frost. The light and dark so
        life-like harmonise with the figure of those  there in the wind,  That when
        I've done tracing their autumn growth, a fragrant smell  issues under my
        wrist.}
}
{
}

\Verse{53}
{
    \zA{}
}
{
    \eA{Do you not mark how they resemble those, by the east hedge, which you
        leisurely pluck? Upon the screens their image I affix to solace me for those
        of the  ninth moon. "Asking the chrysanthemums," by the "Hsiao Hsiang
        consort." Your heart, in autumn, I would like to read, but know it no one
        could! While humming with my arms behind my back, on the east hedge I rap.
        So peerless and unique are ye that who is meet with you to stay? Why are you
        of all flowers the only ones to burst the last in bloom? Why in such silence
        plunge the garden dew and the frost in the hall? When wild geese homeward
        fly and crickets sicken, do you think of me? Do not tell me that in the
        world none of you grow with power of  speech? But if ye fathom what I say,
        why not converse with me a while?}
}
{
}

\Verse{54}
{
    \zA{}
}
{
    \eA{"Pinning the chrysanthemums in the hair," by the "Visitor under the banana
        trees." I put some in a vase, and plant some by the hedge, so day by day I
        have ample to do. I pluck them, yet don't fancy they are meant for girls to
        pin before  the glass in their coiffure. My mania for these flowers is just
        as keen as was that of the squire,  who once lived in Ch'ang An. I rave as
        much for them as raved Mr. P'eng Tsê, when he was under the  effects of
        wine. Cold is the short hair on his temples and moistened with dew, which
        on  it dripped from the three paths. His flaxen turban is suffused with the
        sweet fragrance of the autumn  frost in the ninth moon. That strong weakness
        of mine to pin them in my hair is viewed with  sneers by my contemporaries.}
}
{
}

\Verse{55}
{
    \zA{}
}
{
    \eA{They clap their hands, but they are free to laugh at me by the  roadside as
        much us e'er they list. "The shadow of the chrysanthemums," by the "Old
        Friend of the hall reclining on the russet clouds." In layers upon layers
        their autumn splendour grows and e'er thick and  thicker. I make off
        furtively, and stealthily transplant them from the three  crossways. The
        distant lamp, inside the window-frame, depicts their shade both  far and
        near. The hedge riddles the moon's rays, like unto a sieve, but the flowers
        stop the holes. As their reflection cold and fragrant tarries here, their
        soul must  too abide. The dew-dry spot beneath the flowers is so like them
        that what is  said  of dreams is trash.}
}
{
}

\Verse{56}
{
    \zA{}
}
{
    \eA{Their precious shadows, full of subtle scent, are trodden down to  pieces
        here and there. Could any one with eyes half closed from drinking, not
        mistake the  shadow for the flowers. "Dreaming of chrysanthemums," by the
        "Hsiao Hsiang consort." What vivid dreams arise as I dose by the hedge
        amidst those autumn  scenes! Whether clouds bear me company or the moon be
        my mate, I can't  discern. In fairyland I soar, not that I would become a
        butterfly like Chang. So long I for my old friend T'ao, the magistrate, that
        I again seek  him. In a sound sleep I fell; but so soon as the wild geese
        cried, they  broke my rest. The chirp of the cicadas gave me such a start
        that I bear them a  grudge. My secret wrongs to whom can I go and divulge,
        when I wake up from  sleep?}
}
{
}

\Verse{57}
{
    \zA{}
}
{
    \eA{The faded flowers and the cold mist make my feelings of anguish know  no
        bounds. "Fading of the chrysanthemums," by the "Visitor under the banana
        trees." The dew congeals; the frost waxes in weight; and gradually dwindles
        their bloom. After the feast, with the flower show, follows the season of
        the  'little snow.' The stalks retain still some redundant smell, but the
        flowers' golden  tinge is faint. The stems do not bear sign of even one
        whole leaf; their verdure is  all past. Naught but the chirp of crickets
        strikes my ear, while the moon  shines  on half my bed. Near the cold
        clouds, distant a thousand li, a flock of wild geese  slowly fly. When
        autumn breaks again next year, I feel certain that we will meet  once more.
        We part, but only for a time, so don't let us indulge in anxious  thoughts.}
}
{
}

\Verse{58}
{
    \zA{}
}
{
    \eA{Each stanza they read they praised; and they heaped upon each other
        incessant eulogiums. "Let me now criticise them; I'll do so with all
        fairness!" Li Wan smiled. "As I glance over the page," she said, "I find
        that each of you has some distinct admirable sentiments; but in order to be
        impartial in my criticism to-day, I must concede the first place to:
        'Singing the chrysanthemums;' the second to: 'Asking the chrysanthemums;'
        and the third to: 'Dreaming of chrysanthemums.' The original nature of the
        themes makes the verses full of originality, and their conception still more
        original. But we must allow to the 'Hsiao Hsiang consort' the credit of
        being the best;}
}
{
}

\Verse{59}
{
    \zA{}
}
{
    \eA{next in order following: 'Pinning chrysanthemums in the hair,' 'Facing the
        chrysanthemums,' 'Putting the chrysanthemums, in vases,' 'Drawing the
        chrysanthemums,' and 'Longing for chrysanthemums,' as second best." This
        decision filled Pao-yü with intense gratification. Clapping his hands,
        "Quite right! it's most just," he shouted. "My verses are worth nothing!"
        Tai-yü remarked. "Their fault, after all, is that they are a little too
        minutely subtile." "They are subtile but good," Li Wan rejoined; "for
        there's no artificialness or stiffness about them." "According to my views,"
        Tai-yü observed, "the best line is:  "'When cold holdeth the park and the
        sun's rays do slant, I long and  yearn for you, old friends.' "The metonomy:
        "'I fling my book aside and turn my gaze upon a twig of autumn.'}
}
{
}

\Verse{60}
{
    \zA{}
}
{
    \eA{is already admirable! She has dealt so exhaustively with 'putting
        chrysanthemums in a vase' that she has left nothing unsaid that could be
        said, and has had in consequence to turn her thought back and consider the
        time anterior to their being plucked and placed in vases. Her sentiments are
        profound!" "What you say is certainly so," explained Li Wan smiling; "but
        that line of yours:  "'Some scent I hold by the side of my mouth,….' "beats
        that." "After all," said T'an Ch'un, "we must admit that there's depth of
        thought in those of the 'Princess of Heng Wu' with:  "'…in autumn all trace
        of you is gone;' "and  "'…my dreams then know something of you!' "They
        really make the meaning implied by the words 'long for' stand out clearly."}
}
{
}

\Verse{61}
{
    \zA{}
}
{
    \eA{"Those passages of yours:  "'Cold is the short hair on his temples and
        moistened….' "and  "'His flaxen turban is suffused with the sweet
        fragrance….;'" laughingly observed Puo-ch'ai, "likewise bring out the idea
        of 'pinning the chrysanthemums in the hair' so thoroughly that one couldn't
        get a loop hole for fault-finding." Hsiang-yün then smiled. "'…who is meet
        with you to stay'" she said, "and  "'…burst the last in bloom.' "are
        questions so straight to the point set to the chrysanthemums, that they are
        quite at a loss what answer to give." "Were what you say:  "'I sit with head
        uncovered….' "and  "'…clasping my knees, I hum my lays….'}
}
{
}

\Verse{62}
{
    \zA{}
}
{
    \eA{"as if you couldn't, in fact, tear yourself away for even a moment from
        them," Li Wan laughed, "to come to the knowledge of the chrysanthemums, why,
        they would certainly be sick and tired of you." This joke made every one
        laugh. "I'm last again!" smiled Pao-yü. "Is it likely that:  "'Who plants
        the flowers?…. …in autumn where do they go? With sandals  waxed I come from
        distant shores;…. …and as on this cold day I can't  exhaust my song;….' "do
        not all forsooth amount to searching for chrysanthemums? And that  "'Last
        night they got a shower…. And this morn … bedecked with frost,' "don't both
        bear on planting them? But unfortunately they can't come up to these lines:
        "'Some scent I hold by the side of my mouth and turning to the moon I  sing
        my sentiments.'}
}
{
}

\Verse{63}
{
    \zA{}
}
{
    \eA{'In their pure and cool fragrance, clasping my  knees I hum my lays.'
        '…short hair on his temples….' 'His flaxen  turban…. …golden tinge is faint.
        …verdure is all past. …in autumn …  all trace of you is gone. …my dreams
        then know something of you.' "But to-morrow," he proceeded, "if I have got
        nothing to do, I'll write twelve stanzas my self." "Yours are also good," Li
        Wan pursued, "the only thing is that they aren't as full of original
        conception as those other lines, that's all." But after a few further
        criticisms, they asked for some more warm crabs; and, helping themselves, as
        soon as they were brought, from the large circular table, they regaled
        themselves for a time.}
}
{
}

\Verse{64}
{
    \zA{}
}
{
    \eA{"With the crabs to-day in one's hand and the olea before one's eyes, one
        cannot help inditing verses," Pao-yü smiled. "I've already thought of a few;
        but will any of you again have the pluck to devise any?" With this
        challenge, he there and then hastily washed his hands and picking up a pen
        he wrote out what, his companions found on perusal, to run in this strain:
        When in my hands I clasp a crab what most enchants my heart is the  cassia's
        cool shade. While I pour vinegar and ground ginger, I feel from joy as if I
        would  go mad. With so much gluttony the prince's grandson eats his crabs
        that he  should have some wine.}
}
{
}

\Verse{65}
{
    \zA{}
}
{
    \eA{The side-walking young gentleman has no intestines in his frame at  all. I
        lose sight in my greediness that in my stomach cold accumulates. To my
        fingers a strong smell doth adhere and though I wash them yet  the smell
        clings fast. The main secret of this is that men in this world make much of
        food. The P'o Spirit has laughed at them that all their lives they only
        seek  to eat. "I could readily compose a hundred stanzas with such verses in
        no time," Tai-yü observed with a sarcastic smile. "Your mental energies are
        now long ago exhausted," Pao-yü rejoined laughingly, "and instead of
        confessing your inability to devise any, you still go on heaping invective
        upon people!" Tai-yü, upon catching this insinuation, made no reply of any
        kind;}
}
{
}

\Verse{66}
{
    \zA{}
}
{
    \eA{but slightly raising her head she hummed something to herself for a while,
        and then taking up a pen she completed a whole stanza with a few dashes. The
        company then read her lines. They consisted of—  E'en after death, their
        armour and their lengthy spears are never  cast  away. So nice they look,
        piled in the plate, that first to taste them I'd  fain be. In every pair of
        legs they have, the crabs are full of tender  jade-like meat. Each piece of
        ruddy fat, which in their shell bumps up, emits a  fragrant smell. Besides
        much meat, they have a greater relish for me still, eight  feet  as well.
        Who bids me drink a thousand cups of wine in order to enhance my joy?}
}
{
}

\Verse{67}
{
    \zA{}
}
{
    \eA{What time I can behold their luscious food, with the fine season doth
        accord  When cassias wave with fragrance pure, and the chrysanthemums are
        decked with frost. Pao-yü had just finished conning it over and was
        beginning to sing its praise, when Tai-yü, with one snatch, tore it to
        pieces and bade a servant go and burn it. "As my compositions can't come up
        to yours," she then observed, "I'll burn it. Yours is capital, much better
        than the lines you wrote a little time back on the chrysanthemums, so keep
        it for the benefit of others." "I've likewise succeeded, after much effort,
        in putting together a stanza," Pao-ch'ai laughingly remarked. "It cannot, of
        course, be worth much, but I'll put it down for fun's sake." As she spoke,
        she too wrote down her lines.}
}
{
}

\Verse{68}
{
    \zA{}
}
{
    \eA{When they came to look at them, they read—  On this bright beauteous day, I
        bask in the dryandra shade, with a  cup  in my hand. When I was at Ch'ang
        An, with drivelling mouth, I longed for the  ninth  day of the ninth moon.
        The road stretches before their very eyes, but they can't tell  between
        straight and transverse. Under their shells in spring and autumn only reigns
        a vacuum, yellow  and black. At this point, they felt unable to refrain from
        shouting: "Excellent!" "She abuses in fine style!" Pao-yü shouted. "But my
        lines should also be committed to the flames." The company thereupon scanned
        the remainder of the stanza, which was couched in this wise:  When all the
        stock of wine is gone, chrysanthemums then use to scour  away the smell.}
}
{
}

\Verse{69}
{
    \zA{}
}
{
    \eA{So as to counteract their properties of gath'ring cold, fresh ginger  you
        should take. Alas! now that they have been dropped into the boiling pot,
        what good  do they derive? About the moonlit river banks there but remains
        the fragrant aroma of  corn. At the close of their perusal, they with one
        voice, explained that this was a first-rate song on crab-eating; that minor
        themes of this kind should really conceal lofty thoughts, before they could
        be held to be of any great merit, and that the only thing was that it
        chaffed people rather too virulently. But while they were engaged in
        conversation, P'ing Erh was again seen coming into the garden. What she
        wanted is not, however, yet known; so, reader, peruse the details given in
        the subsequent chapter.}
}
{
}
